\documentclass[12pt]{amsart}

\usepackage{amsmath,amssymb,amsthm}
\usepackage{mathtools}
\usepackage{booktabs}
\usepackage[margin=1.1in]{geometry}
\usepackage[
  pdftitle={Prime magnitudes and nonvanishing for a nonmultiplicative indefinite theta series over Q(sqrt(5))},
  pdfauthor={Xiang Huang}
]{hyperref}

\newtheorem{theorem}{Theorem}[section]
\newtheorem{proposition}[theorem]{Proposition}
\newtheorem{lemma}[theorem]{Lemma}
\newtheorem{corollary}[theorem]{Corollary}
\newtheorem{conjecture}[theorem]{Conjecture}
\theoremstyle{definition}
\newtheorem{definition}[theorem]{Definition}
\newtheorem{example}[theorem]{Example}
\theoremstyle{remark}
\newtheorem{remark}[theorem]{Remark}

\DeclareMathOperator{\N}{N}
\DeclareMathOperator{\Tr}{Tr}
\newcommand{\ZZ}{\mathbb{Z}}
\newcommand{\QQ}{\mathbb{Q}}
\newcommand{\RR}{\mathbb{R}}
\newcommand{\FF}{\mathbb{F}}
\newcommand{\OK}{\mathcal{O}_K}
\newcommand{\eps}{\varepsilon}
\newcommand{\ph}{\varphi}
\newcommand{\Acone}{\mathcal{A}}
\newcommand{\Dcone}{\mathcal{D}}
\newcommand{\Sector}{\mathcal{S}}

\title[Prime magnitudes of a nonmultiplicative theta series]
{Prime magnitudes and nonvanishing for a nonmultiplicative\\
indefinite theta series over $\QQ(\sqrt5)$}
\author{Xiang Huang}
\date{}

\begin{document}

\begin{abstract}
We study an independently defined norm-supported cone series
$B(q)=\sum_{N\geq0}B_Nq^N$ over $\QQ(\sqrt5)$.  Its coefficients are
signed counts of points in two quadrants of an affine lattice coset,
and the norm parametrization is
$-\N(\beta(k,r))=10E(k,r)+1$.  The sequence is not multiplicative:
with $a(M)=B_{(M-1)/10}$ one has
$a(11)a(31)=-2\neq3=a(341)$.  Nevertheless its values at rational
primes are rigid.  For every prime $p\equiv1\pmod{10}$ we construct a
three-valued ideal invariant
\[
  \iota(\mathfrak p)=\delta(\pi)-\lambda(\pi)\in\ZZ/3\ZZ,
\]
where $\delta$ is an archimedean sector index and $\lambda$ is the
discrete logarithm of $\pi$ in the group
$(\OK/2\OK)^\times\simeq\FF_4^\times$.  The invariant combines an
archimedean sector label with a congruence label; we do not identify it
with a finite ray-class character.  We prove a reflection identity, an
exact one-atom lemma, and a sign-coherence lemma,
yielding
\[
  B_{(p-1)/10}\in\{-2,-1,1,2\},
  \qquad
  |B_{(p-1)/10}|=2\iff\iota(\mathfrak p)=1.
\]
In particular the prime coefficients never vanish.  A specialization
of the Hecke--Mitsui prime number theorem to twelve joint
sector--congruence cells gives
\[
 \#\{p\leq X:p\equiv1\pmod{10},\ |B_{(p-1)/10}|=2\}
 =\frac1{12}\operatorname{Li}(X)+o(X/\log X)
\]
and the corresponding count for magnitude $1$ is
$\frac16\operatorname{Li}(X)+o(X/\log X)$.  Thus the relative
densities are $1/3$ and $2/3$.  Lean~4 checks the complete algebraic,
geometric, and finite prime-classification argument: generator
existence, exact atom counts, sign coherence, and nonvanishing.  The
specialized Hecke--Mitsui prime-element asymptotic is the
external analytic input and is not formalized in this public Lean package.
The constants $1/12$, $1/6$, $1/3$, and $2/3$ are deduced in the manuscript
from the six selected-cell asymptotics and the two-ideal bookkeeping below;
no density theorem is part of the Comparator-selected Lean surface.
\end{abstract}

\maketitle

%----------------------------------------------------------------------
\section{Introduction}\label{sec:intro}
%----------------------------------------------------------------------

The Andrews--Dyson--Hickerson and Cohen identities connect
real-quadratic norm sums with Maass waveforms and multiplicative
coefficients \cite{jADH88,jCo88,jCo95}.  In several standard examples,
the defining lattice or coset is stable under the fundamental totally
positive unit, so the element sum descends to a Hecke-character sum
over ideals and multiplicativity follows.  This mechanism occurs in
the original construction over $\QQ(\sqrt6)$
\cite{jADH88,jCo88}, the $\QQ(\sqrt2)$ example of
Corson--Favero--Liesinger--Zubairy \cite{jCFLZ04}, Lovejoy's
overpartition series \cite{jLo04}, the multiplicative
$q$-hypergeometric families of Bringmann--Kane \cite{jBrKa11}, and
the double sums of Lovejoy--Osburn \cite{jLO13,jLoOs15}.  The series
considered here has the same broad norm-supported shape, but its
defining coset is \emph{not} unit-stable: it does not descend to a
single Hecke-character sum over ideals, and its coefficients are not
multiplicative.  The resulting nontrivial unit phase is the structural
feature studied below; no priority claim is needed for our arguments.

Let
\[
  K=\QQ(\sqrt5),
  \qquad
  \OK=\ZZ[\ph],
  \qquad
  \ph=\frac{1+\sqrt5}{2},
  \qquad
  \eps=\ph^2=\ph+1.
\]
The unit $\ph$ has norm $-1$, while $\eps$ is the fundamental totally
positive unit.  We define an affine coset $L\subset\OK$, a bijection
$\beta:\ZZ^2\to L$, and a cone-difference series $B$.  The paper has
three main goals:
\begin{enumerate}
\item prove directly that the classical multiplicative mechanism is
unavailable and give an explicit multiplicativity counterexample;
\item classify the magnitudes of all coefficients at primes
$p\equiv1\pmod{10}$ and prove nonvanishing;
\item prove the absolute prime-counting asymptotics, and hence the
relative densities $1/3$ and $2/3$, using joint
sector-and-congruence equidistribution of prime elements.
\end{enumerate}

The definition of the finite invariant combines an archimedean sector
index with a congruence index.  We do not prove that it factors through
a finite ray-class quotient, and we do not use such a factorization.
The density argument is instead formulated in terms of a joint
sector-and-congruence prime-element asymptotic of Hecke--Mitsui type
\cite{jHe20,jMi56,oKai22}.

Here is the main result.  The invariant $\iota$ is defined in
Section~\ref{sec:iota}.

\begin{theorem}[Prime classification and density]\label{thm:main}
For a rational prime $p\equiv1\pmod{10}$, put
\[
  a(p):=B_{(p-1)/10}.
\]
(These are the only primes that occur: the norm identity represents
exactly the values $10N+1$, so $p=5$ and the inert primes
$p\equiv\pm3\pmod{10}$ never appear, and the primes
$p\equiv9\pmod{10}$ split in $\QQ(\sqrt5)$ but are not of the form
$10N+1$ and hence never enter the indexing.)
Then:
\begin{enumerate}
\item $a(p)\in\{-2,-1,1,2\}$, so $a(p)\neq0$;
\item for either prime ideal $\mathfrak p\mid p$,
\[
  |a(p)|=2\quad\Longleftrightarrow\quad
  \iota(\mathfrak p)=1;
\]
\item if $P_2(X)$ and $P_1(X)$ count the primes $p\leq X$,
$p\equiv1\pmod{10}$, for which $|a(p)|$ is respectively $2$ and $1$,
then
\[
  P_2(X)=\frac1{12}\operatorname{Li}(X)
    +o\!\left(\frac{X}{\log X}\right),
  \qquad
  P_1(X)=\frac16\operatorname{Li}(X)
    +o\!\left(\frac{X}{\log X}\right).
\]
Consequently the relative densities are $1/3$ and $2/3$.
\end{enumerate}
\end{theorem}

The series was suggested by calculations around Chan's dissection of
$\Theta_{10}$ \cite{jCh10,oCh11}.  The present paper deliberately
starts from the independent definition of $B$ in
Definition~\ref{def:B}.  A row-model factorization motivated that
definition and has been checked through a finite coefficient range,
but we do not claim here a source-faithful theorem identifying Chan's
missing kernel with $B$; no result below depends on such an
identification.

Section~\ref{sec:algebra} gives the lattice and norm formulas.
Section~\ref{sec:units} studies the unit action and proves
nonmultiplicativity.  Sections~\ref{sec:sectors} and \ref{sec:iota}
construct the sector--residue invariant.  Section~\ref{sec:prime}
proves the one-atom and sign-coherence lemmas and derives prime
nonvanishing.  Section~\ref{sec:density} derives the density statement
from its explicitly isolated analytic input.
Section~\ref{sec:verification} records computational and formalization
scope.

%----------------------------------------------------------------------
\section{The lattice, cones, and norm parametrization}
\label{sec:algebra}
%----------------------------------------------------------------------

\subsection{Basic arithmetic of the quadratic field}

The conjugate of $\ph$ is $\bar\ph=1-\ph$.  For
$\alpha=a+b\ph\in\OK$,
\[
  \bar\alpha=(a+b)-b\ph,
  \qquad
  \Tr(\alpha)=2a+b,
  \qquad
  \N(\alpha)=a^2+ab-b^2.
\]
The two real embeddings are denoted by
\[
  \sigma_1(a+b\ph)=a+b\ph,
  \qquad
  \sigma_2(a+b\ph)=a+b\bar\ph.
\]

\begin{lemma}[Class number one]\label{lem:class-number-one}
The ordinary and narrow class numbers of $K$ are both one.
\end{lemma}

\begin{proof}
The discriminant of $K$ is $5$.  Minkowski's bound in degree two says
that every ideal class contains an integral ideal of norm at most
\[
  \frac{2!}{2^2}\sqrt5=\frac{\sqrt5}{2}<2.
\]
The norm of a nonzero integral ideal is a positive integer, so the
representative has norm $1$ and is $\OK$.  Hence the ordinary class
number is one.  Since $\ph$ has norm $-1$, every principal ideal has a
generator with either prescribed sign pattern, and the narrow class
number is also one.
\end{proof}

\begin{lemma}[Splitting of the relevant primes]\label{lem:splitting}
Every rational prime $p\equiv1\pmod{10}$ splits as
$p\OK=\mathfrak p\bar{\mathfrak p}$ with two distinct principal prime
ideals of norm $p$.
\end{lemma}

\begin{proof}
Such a prime is odd and is not $5$.  Quadratic reciprocity gives
\[
  \left(\frac{5}{p}\right)=\left(\frac{p}{5}\right)=1
\]
because $p\equiv1\pmod5$.  Thus $p$ splits in $K$.  Principality
follows from Lemma~\ref{lem:class-number-one}.
\end{proof}

\subsection{The affine coset and the exponent}

\begin{definition}\label{def:L}
Define
\[
  L:=\{a+b\ph\in\OK:b-3a\equiv1\pmod{10}\}.
\]
This is a translate of the index-$10$ sublattice
\[
  L_0:=\{a+b\ph:b-3a\equiv0\pmod{10}\}.
\]
\end{definition}

\begin{definition}\label{def:beta-E}
For $k,r\in\ZZ$, define
\[
  \beta(k,r):=(r-2k)+(4k+3r+1)\ph
\]
and
\[
  H(k,r):=4k^2+2k+r^2+(6k+1)r,
  \qquad
  E(k,r):=\frac{H(k,r)}2.
\]
Define the two cones
\[
  \Acone:=\{(k,r)\in\ZZ^2:k\geq0,\ r\geq0\},
  \qquad
  \Dcone:=\{(k,r)\in\ZZ^2:k<0,\ r<0\}.
\]
\end{definition}

\begin{lemma}[Integrality and cone nonnegativity]\label{lem:E}
For every $k,r\in\ZZ$, the integer $H(k,r)$ is even, so
$E(k,r)\in\ZZ$.  Moreover $E(k,r)\geq0$ on
$\Acone\cup\Dcone$.
\end{lemma}

\begin{proof}
We have
\[
  H(k,r)=r(r+1)+2k(2k+1)+6kr,
\]
a sum of two even products and an even multiple.  Hence $H$ is even.
If $(k,r)\in\Acone$, then
\[
  E(k,r)=2k^2+k+3kr+\frac{r(r+1)}2\geq0.
\]
If $(k,r)\in\Dcone$, write $k=-m-1$ and $r=-s-1$ with
$m,s\geq0$.  Then
\[
  E(-m-1,-s-1)
  =2m^2+6m+4+3ms+3s+\frac{s(s+1)}2\geq0.
\]
No nonnegativity assertion is made outside the two cones.
\end{proof}

\begin{proposition}[Bijection and norm identity]\label{prop:beta-norm}
The map $\beta:\ZZ^2\to L$ is a bijection.  Its inverse is
\[
  k=\frac{b-3a-1}{10},
  \qquad
  r=a+2k
\]
for $a+b\ph\in L$.  Furthermore,
\[
  -\N(\beta(k,r))=10E(k,r)+1.
\]
\end{proposition}

\begin{proof}
For $\beta(k,r)=a+b\ph$,
\[
  b-3a-1
  =(4k+3r+1)-3(r-2k)-1=10k,
\]
so $\beta(k,r)\in L$.  The displayed inverse formulas are integral for
an element of $L$ and recover $k,r$, proving bijectivity.  Substituting
$a=r-2k$ and $b=4k+3r+1$ into
$\N(a+b\ph)=a^2+ab-b^2$ gives
\[
  \N(\beta(k,r))=-10E(k,r)-1.
\]
\end{proof}

\begin{definition}[The cone series]\label{def:B}
For $N\geq0$, define
\[
  B_N:=
  \sum_{\substack{(k,r)\in\Dcone\\E(k,r)=N}}(-1)^r
  -
  \sum_{\substack{(k,r)\in\Acone\\E(k,r)=N}}(-1)^r,
  \qquad
  B(q):=\sum_{N\geq0}B_Nq^N.
\]
For later use, define the contribution of a single atom by
\[
  w(k,r):=
  \begin{cases}
    -(-1)^r,&(k,r)\in\Acone,\\
    (+1)(-1)^r,&(k,r)\in\Dcone,\\
    0,&\text{otherwise}.
  \end{cases}
\]
\end{definition}

\begin{lemma}[Finiteness of coefficients]\label{lem:finite}
For every $N\geq0$, both sums in Definition~\ref{def:B} are finite.
\end{lemma}

\begin{proof}
On $\Acone$ the expression in Lemma~\ref{lem:E} gives
$2k^2+k\leq N$ and $r(r+1)/2\leq N$.  Hence $k$ and $r$ are bounded.
On $\Dcone$, the formula with $k=-m-1$, $r=-s-1$ gives
$2m^2\leq N$ and $s(s+1)/2\leq N$, so $m$ and $s$ are bounded.
\end{proof}

\begin{corollary}[Norm support]\label{cor:norm-support}
If $B_N\neq0$, then $10N+1$ is the norm of an algebraic integer of
$K$.
\end{corollary}

\begin{proof}
A nonzero coefficient contains an atom $\beta(k,r)$ of norm
$-(10N+1)$.  Multiplication by $\ph$, whose norm is $-1$, changes this
to norm $10N+1$.
\end{proof}

\begin{remark}[Provenance]
The row-model factorization that motivated $B$ was verified
coefficient-by-coefficient through $N=2048$.  This finite check is not
a theorem identifying the original missing kernel with $B$, and the
proof of the present paper does not use it.
\end{remark}

%----------------------------------------------------------------------
\section{Unit action and nonmultiplicativity}\label{sec:units}
%----------------------------------------------------------------------

The congruence defining $L$ has one component modulo $2$ and one
component modulo the ramified prime $\sqrt5$.

\begin{lemma}[Two congruence descriptions of $L$]\label{lem:L-CRT}
For $\alpha=a+b\ph\in\OK$, the following are equivalent:
\begin{enumerate}
\item $\alpha\in L$;
\item
\[
  \alpha\equiv3\pmod{\sqrt5\OK}
  \quad\text{and}\quad
  \alpha\bmod2\in\{1,\ph\}\subset\FF_4^\times.
\]
\end{enumerate}
\end{lemma}

\begin{proof}
Modulo $\sqrt5$ one has $\ph\equiv3\pmod5$, so
\[
  \alpha\equiv a+3b\pmod{\sqrt5}.
\]
If $c=b-3a$, then $a+3b\equiv3c\pmod5$.  Thus
$c\equiv1\pmod5$ if and only if
$\alpha\equiv3\pmod{\sqrt5}$.  Modulo $2$, the condition
$b-3a\equiv1$ is $a+b\equiv1$, which says that exactly one of $a,b$
is odd.  These are precisely the residues $1$ and $\ph$ in
$\OK/2\OK\simeq\FF_4$.  Combining the mod-$2$ and mod-$5$ conditions
by the Chinese remainder theorem gives the mod-$10$ condition.
\end{proof}

\begin{proposition}[Unit instability and stabilizer]\label{prop:unit}
One has
\[
  \eps L\cap L=\varnothing.
\]
Moreover, for $j\in\ZZ$,
\[
  \eps^jL=L\quad\Longleftrightarrow\quad6\mid j.
\]
\end{proposition}

\begin{proof}
Modulo $\sqrt5$, one has
\[
  \eps=\ph^2\equiv3^2\equiv-1\pmod5.
\]
Thus $\eps$ has order $2$ in the $\sqrt5$ component.  Since every
element of $L$ is congruent to $3$ modulo $\sqrt5$, multiplication by
$\eps$ changes that residue to $-3\equiv2$, proving
$\eps L\cap L=\varnothing$.  More generally, preservation of the
$\sqrt5$ condition requires $j$ even.

Modulo $2$, the group $\FF_4^\times$ is cyclic of order $3$ and
$\eps\equiv\ph+1$ is a generator.  The two residues allowed by
Lemma~\ref{lem:L-CRT} have discrete logarithms $0$ and $2$.  Translation
of the set $\{0,2\}$ by $j$ preserves it if and only if
$j\equiv0\pmod3$.  Hence $\eps^jL=L$ exactly when $j$ is divisible by
both $2$ and $3$, equivalently by $6$.
\end{proof}

The proposition explains why the standard one-coset descent to ideals
is unavailable.  It does not by itself prove that the coefficient
sequence is nonmultiplicative; the following finite witness does.

\begin{theorem}[Nonmultiplicativity]\label{thm:nonmult}
For $M\equiv1\pmod{10}$ define
\[
  a(M):=B_{(M-1)/10}.
\]
Then
\[
  a(11)=1,
  \qquad
  a(31)=-2,
  \qquad
  a(341)=3,
\]
and therefore
\[
  a(11)a(31)=-2\neq3=a(11\cdot31).
\]
\end{theorem}

\begin{proof}
The complete atom lists are
\begin{center}
\begin{tabular}{ccccc}
\toprule
$N$ & cone & $(k,r)$ & $(-1)^r$ & contribution to $B_N$\\
\midrule
$1$ & $\Acone$ & $(0,1)$ & $-1$ & $+1$\\
\midrule
$3$ & $\Acone$ & $(0,2)$ & $+1$ & $-1$\\
    & $\Acone$ & $(1,0)$ & $+1$ & $-1$\\
\midrule
$34$ & $\Acone$ & $(2,3)$ & $-1$ & $+1$\\
     & $\Dcone$ & $(-1,-6)$ & $+1$ & $+1$\\
     & $\Dcone$ & $(-3,-2)$ & $+1$ & $+1$\\
\bottomrule
\end{tabular}
\end{center}
For the A-cone, the inequalities in the proof of
Lemma~\ref{lem:finite} bound $k$ and $r$ before substitution.  For the
D-cone, use $k=-m-1$, $r=-s-1$ and the corresponding bounds.  Checking
those finite ranges gives exactly the displayed solutions.  Summing
the contributions gives the three values and the claimed inequality.
\end{proof}

%----------------------------------------------------------------------
\section{A half-open sector and the cone criterion}\label{sec:sectors}
%----------------------------------------------------------------------

For a nonzero $\alpha\in K$ of negative norm, define
\[
  R(\alpha):=\left|\frac{\sigma_1(\alpha)}{\sigma_2(\alpha)}\right|,
  \qquad
  \rho:=\ph^4.
\]
Both embeddings of $\eps$ are positive and
\[
  R(\eps\alpha)=\rho R(\alpha).
\]
We use the half-open fundamental sector
\[
  \Sector:=\{\alpha:\sigma_1(\alpha)>0>\sigma_2(\alpha),\
                  1\leq R(\alpha)<\rho\}.
\]
The opposite sector is $-\Sector$.

\begin{lemma}[Prime elements do not lie on sector boundaries]
\label{lem:no-boundary}
Let $p\neq5$ be a rational prime and let
$\alpha\in\OK$ satisfy $\N(\alpha)=-p$.  Then
$R(\alpha)\neq\rho^j$ for every $j\in\ZZ$.
\end{lemma}

\begin{proof}
If $R(\alpha)=\rho^j$, then
$\gamma=\eps^{-j}\alpha$ has $R(\gamma)=1$.  After replacing
$\gamma$ by $-\gamma$ if necessary, its embeddings have opposite signs
and equal absolute value, so $\Tr(\gamma)=0$.  Writing
$\gamma=a+b\ph$, the equation $2a+b=0$ implies that $b$ is even and
$\gamma$ is an integral multiple of $\sqrt5$.  Hence
$\N(\gamma)=-5m^2$, which cannot equal $-p$ for $p\neq5$.
\end{proof}

\begin{definition}[Sector index]\label{def:delta}
Let $\alpha\in\OK$ have norm $-p$ for a rational prime $p\neq5$.
By Lemma~\ref{lem:no-boundary}, there is a unique
$\delta(\alpha)\in\ZZ$ such that
\[
  \rho^{\delta(\alpha)}<R(\alpha)<
  \rho^{\delta(\alpha)+1}.
\]
Equivalently, $\eps^{-\delta(\alpha)}\alpha$ has ratio strictly between
$1$ and $\rho$.  We use the reduction of $\delta(\alpha)$ modulo $3$
when forming $\iota$.
\end{definition}

The half-open convention fixes the boundary allocation in the
general definition, and Lemma~\ref{lem:no-boundary} shows that no
relevant prime element is affected by it.

\begin{lemma}[Behavior of the sector index]\label{lem:delta}
For every $m\in\ZZ$,
\[
  \delta(\eps^m\alpha)=\delta(\alpha)+m,
  \qquad
  \delta(-\alpha)=\delta(\alpha),
\]
and
\[
  \delta(\bar\alpha)=-\delta(\alpha)-1.
\]
\end{lemma}

\begin{proof}
The first two identities follow from
$R(\eps^m\alpha)=\rho^mR(\alpha)$ and $R(-\alpha)=R(\alpha)$.  For
conjugation, $R(\bar\alpha)=R(\alpha)^{-1}$.  If
$\rho^d<R(\alpha)<\rho^{d+1}$, then
\[
  \rho^{-d-1}<R(\bar\alpha)<\rho^{-d},
\]
which gives the last identity.
\end{proof}

The next proposition is the exact link between the geometric sector
and the two cones in Definition~\ref{def:beta-E}.

\begin{proposition}[Cone recognition]\label{prop:cone-recognition}
Let $\alpha=\beta(k,r)$.  Then
\begin{align*}
(k,r)\in\Acone
&\quad\Longleftrightarrow\quad
\sigma_1(\alpha)>0>\sigma_2(\alpha)
\ \text{and}\ 1<R(\alpha)<\rho,\\
(k,r)\in\Dcone
&\quad\Longleftrightarrow\quad
\sigma_1(\alpha)<0<\sigma_2(\alpha)
\ \text{and}\ 1<R(\alpha)<\rho.
\end{align*}
\end{proposition}

\begin{proof}
Write $\alpha=a+b\ph$.  Direct substitution gives
\[
  \Tr(\alpha)=2a+b=5r+1
\]
and, using $1+\rho=3\ph^2$ and
$\ph+\rho\bar\ph=-\ph^2$,
\[
  \sigma_1(\alpha)+\rho\sigma_2(\alpha)
  =\ph^2(3a-b)=-\ph^2(10k+1).
\]
Suppose first that $\sigma_1(\alpha)>0>\sigma_2(\alpha)$.  Then
$R(\alpha)>1$ is equivalent to
$\Tr(\alpha)>0$, hence to $5r+1>0$, which for integral $r$ is
$r\geq0$.  Likewise $R(\alpha)<\rho$ is equivalent to
$\sigma_1(\alpha)+\rho\sigma_2(\alpha)<0$, hence to
$10k+1>0$, which is $k\geq0$.

Conversely, if $k,r\geq0$, then
\[
  \beta(k,r)=\ph+2k\sqrt5+r(1+3\ph).
\]
Each summand has positive first embedding and negative second
embedding.  The two displayed linear identities then give
$1<R(\alpha)<\rho$.

For the D-cone apply the same argument to $-\alpha$.  The inequalities
become $-(5r+1)>0$ and $10k+1<0$, equivalent to
$r\leq-1$ and $k\leq-1$.  Conversely, writing
$k=-m-1$, $r=-s-1$ gives
\[
  -\beta(k,r)=(-1+6\ph)+2m\sqrt5+s(1+3\ph),
\]
whose summands again have positive first embedding and negative second
embedding.  This proves the second equivalence.
\end{proof}

%----------------------------------------------------------------------
\section{The sector--residue invariant}\label{sec:iota}
%----------------------------------------------------------------------

Because $2$ is inert in $K$,
\[
  \OK/2\OK\simeq\FF_4,
  \qquad
  \FF_4^\times=\langle\eps\bmod2\rangle.
\]
For $\alpha\not\equiv0\pmod2$, define
\[
  \lambda(\alpha)\in\ZZ/3\ZZ
\]
by
\[
  \alpha\equiv\eps^{\lambda(\alpha)}\pmod2.
\]
Thus
\[
  \lambda(\eps^m\alpha)=\lambda(\alpha)+m,
  \qquad
  \lambda(-\alpha)=\lambda(\alpha).
\]
Conjugation on $\FF_4$ is Frobenius, so
\[
  \lambda(\bar\alpha)=2\lambda(\alpha).
\]

\begin{definition}[The invariant $\iota$]\label{def:iota}
Let $p\equiv1\pmod{10}$ be prime and let
$\mathfrak p\mid p$.  Choose a generator $\pi$ of $\mathfrak p$ with
$\N(\pi)=-p$, and define
\[
  \iota(\mathfrak p)
  :=\delta(\pi)-\lambda(\pi)\in\ZZ/3\ZZ.
\]
\end{definition}

\begin{proposition}[Well-definedness]\label{prop:iota-wd}
The value $\iota(\mathfrak p)$ is independent of the chosen generator
of norm $-p$.
\end{proposition}

\begin{proof}
The units of norm $1$ are $\{\pm\eps^m:m\in\ZZ\}$.  Thus another
generator of norm $-p$ has the form $\pi'=\pm\eps^m\pi$.  By
Lemma~\ref{lem:delta},
\[
  \delta(\pi')=\delta(\pi)+m,
\]
while the discrete logarithm satisfies
\[
  \lambda(\pi')=\lambda(\pi)+m\pmod3.
\]
Their difference is unchanged.
\end{proof}

\begin{theorem}[Reflection identity]\label{thm:reflection}
For conjugate prime ideals,
\[
  \iota(\bar{\mathfrak p})
  \equiv-\iota(\mathfrak p)-1\pmod3.
\]
\end{theorem}

\begin{proof}
Use $\bar\pi$ as a generator of $\bar{\mathfrak p}$.  By
Lemma~\ref{lem:delta},
$\delta(\bar\pi)=-\delta(\pi)-1$.  In $\FF_4$,
$\lambda(\bar\pi)=2\lambda(\pi)$.  Hence
\begin{align*}
\iota(\bar{\mathfrak p})
&=(-\delta(\pi)-1)-2\lambda(\pi)\\
&\equiv-(\delta(\pi)-\lambda(\pi))-1
 \pmod3.
\end{align*}
\end{proof}

\begin{remark}
The invariant has finite codomain, but its definition uses the
archimedean sector label $\delta$ as well as a finite congruence label.
We neither identify it with a ray-class character nor use a
Chebotarev argument.  Section~\ref{sec:density} instead isolates the
joint sector-and-congruence equidistribution statement needed below.
\end{remark}

%----------------------------------------------------------------------
\section{Prime atoms, uniqueness, and sign coherence}\label{sec:prime}
%----------------------------------------------------------------------

We first translate membership in $L$ into the normalized sector data.

\begin{lemma}[Normalization into $L$]\label{lem:normalization}
Let $p\equiv1\pmod{10}$ be prime, let $\mathfrak p\mid p$, and let
$\pi$ generate $\mathfrak p$ with $\N(\pi)=-p$.  Put
\[
  \pi_0:=\eps^{-\delta(\pi)}\pi.
\]
Exactly one of $\pi_0$ and $-\pi_0$ is congruent to $3$ modulo
$\sqrt5$.  Denote that choice by $\alpha$.  Then
\[
  \alpha\in L
  \quad\Longleftrightarrow\quad
  \iota(\mathfrak p)\in\{0,1\}.
\]
When this holds, $\alpha=\beta(k,r)$ for a point in
$\Acone\cup\Dcone$.
\end{lemma}

\begin{proof}
Modulo $\sqrt5$, conjugation is trivial, so
$\N(\pi_0)\equiv\pi_0^2$.  Since
$\N(\pi_0)=-p\equiv-1\equiv4\pmod5$, the residue of $\pi_0$ is
$2$ or $3$.  Exactly one sign therefore gives residue $3$.

The sector normalization gives $1<R(\alpha)<\rho$ by
Lemma~\ref{lem:no-boundary}.  Its discrete logarithm is
\[
  \lambda(\alpha)
  =\lambda(\pi)-\delta(\pi)
  =-\iota(\mathfrak p)\pmod3.
\]
By Lemma~\ref{lem:L-CRT}, membership in $L$ is equivalent to the
already imposed residue $3$ modulo $\sqrt5$ together with
$\lambda(\alpha)\in\{0,2\}$.  The latter condition is equivalent to
$\iota(\mathfrak p)\in\{0,1\}$.  If $\alpha\in L$, write
$\alpha=\beta(k,r)$ by Proposition~\ref{prop:beta-norm}; then
Proposition~\ref{prop:cone-recognition} places $(k,r)$ in the A- or
D-cone according to the sign pattern of the embeddings.
\end{proof}

\begin{theorem}[The unique bad class]\label{thm:bad-class}
A prime ideal $\mathfrak p\mid p$ contributes an atom to
$B_{(p-1)/10}$ if and only if
\[
  \iota(\mathfrak p)\neq2.
\]
\end{theorem}

\begin{proof}
This is exactly Lemma~\ref{lem:normalization}: the contributing classes
are $0$ and $1$, while class $2$ fails the mod-$2$ condition defining
$L$.
\end{proof}

The next result supplies the multiplicity statement needed in the
prime theorem.

\begin{lemma}[One atom per contributing prime ideal]
\label{lem:one-atom}
Let $p\equiv1\pmod{10}$ be prime and let $\mathfrak p\mid p$.  Among
all pairs $(k,r)\in\Acone\cup\Dcone$ for which
\[
  (\beta(k,r))=\mathfrak p,
  \qquad
  \N(\beta(k,r))=-p,
\]
there is at most one.  If $\iota(\mathfrak p)\neq2$, there is exactly
one.
\end{lemma}

\begin{proof}
Existence for $\iota(\mathfrak p)\neq2$ is
Lemma~\ref{lem:normalization}.  For uniqueness, suppose
$\alpha=\beta(k,r)$ and $\alpha'=\beta(k',r')$ are two such atoms.
They generate the same ideal and have the same norm, so
\[
  \alpha'=\pm\eps^m\alpha
\]
for some $m\in\ZZ$.  Proposition~\ref{prop:cone-recognition} gives
\[
  1<R(\alpha)<\rho,
  \qquad
  1<R(\alpha')<\rho.
\]
But
$R(\alpha')=\rho^mR(\alpha)$, and two numbers in the same interval of
multiplicative width $\rho$ can satisfy this only when $m=0$.
Therefore $\alpha'=\pm\alpha$.

Both elements lie in $L$, hence both are congruent to $3$ modulo
$\sqrt5$ by Lemma~\ref{lem:L-CRT}.  Negation changes that residue to
$-3\equiv2$, so $\alpha'=-\alpha$ is impossible.  Thus
$\alpha'=\alpha$, and bijectivity of $\beta$ gives
$(k',r')=(k,r)$.
\end{proof}

When both conjugate prime ideals contribute, their two atoms cannot
cancel.  The proof gives an explicit coordinate relation.

\begin{lemma}[Sign coherence]\label{lem:sign-coherence}
Let $p\equiv1\pmod{10}$ be prime.  Suppose both
$\mathfrak p$ and $\bar{\mathfrak p}$ contribute atoms to
$B_{(p-1)/10}$.  If $\alpha=\beta(k,r)$ is the atom belonging to
$\mathfrak p$, then the atom belonging to $\bar{\mathfrak p}$ is
\[
  \alpha'=-\eps\bar\alpha
          =\beta\left(\frac r2,2k\right).
\]
In particular $r$ is even, the two atom coordinates lie in the same
cone, and their contributions to $B_{(p-1)/10}$ have the same sign.
\end{lemma}

\begin{proof}
If both ideals contribute, Theorem~\ref{thm:bad-class} and the
reflection identity leave only the pair of classes
\[
  \iota(\mathfrak p)=\iota(\bar{\mathfrak p})=1.
\]
For the normalized atom $\alpha$, the sector index is $0$, so
\[
  \lambda(\alpha)=-\iota(\mathfrak p)=2\pmod3.
\]
Thus $\alpha\equiv\eps^2\equiv\ph\pmod2$.  If
$\alpha=a+b\ph$, this says that $a$ is even and $b$ is odd.  Since
$a=r-2k$, the integer $r$ is even.

Set $\alpha'=-\eps\bar\alpha$.  It generates
$\bar{\mathfrak p}$ and has norm $-p$.  Modulo $\sqrt5$, conjugation
fixes the residue, $\eps\equiv-1$, and $\alpha\equiv3$, hence
\[
  \alpha'\equiv-(-1)\cdot3\equiv3\pmod{\sqrt5}.
\]
Modulo $2$, conjugation is Frobenius.  Since
$\alpha\equiv\eps^2$, one has
$\bar\alpha\equiv(\eps^2)^2=\eps$, and therefore
$\alpha'\equiv\eps^2\pmod2$.  Lemma~\ref{lem:L-CRT} gives
$\alpha'\in L$.

Furthermore,
\[
  R(\alpha')=R(\eps\bar\alpha)=\frac{\rho}{R(\alpha)},
\]
which lies strictly between $1$ and $\rho$.  Hence $\alpha'$ is an
atom by Proposition~\ref{prop:cone-recognition}; uniqueness follows
from Lemma~\ref{lem:one-atom}.

It remains to compute its coordinates.  Since
\[
  \bar\alpha=(a+b)-b\ph,
\]
multiplication by $\eps$ and negation give
\[
  -\eps\bar\alpha=-a+(b-a)\ph.
\]
Applying the inverse formulas of Proposition~\ref{prop:beta-norm},
\[
  k'=\frac{(b-a)-3(-a)-1}{10}
     =\frac{2a+b-1}{10}=\frac r2,
  \qquad
  r'=-a+2k'=2k.
\]
If $(k,r)$ is in the A-cone, then $(r/2,2k)$ is in the A-cone; if it
is in the D-cone, then $(r/2,2k)$ is in the D-cone.  Both $r$ and
$r'=2k$ are even.  Thus both internal signs $(-1)^r$ equal $+1$, and
because the cone is the same, Definition~\ref{def:B} assigns the same
outer sign to both atoms.
\end{proof}

\begin{theorem}[Prime magnitudes and nonvanishing]
\label{thm:prime-classification}
For every prime $p\equiv1\pmod{10}$ and either
$\mathfrak p\mid p$,
\[
  B_{(p-1)/10}\in\{-2,-1,1,2\},
\]
and
\[
  \left|B_{(p-1)/10}\right|=2
  \quad\Longleftrightarrow\quad
  \iota(\mathfrak p)=1.
\]
In particular $B_{(p-1)/10}\neq0$.
\end{theorem}

\begin{proof}
Let $i=\iota(\mathfrak p)$.  By Theorem~\ref{thm:reflection}, the
ordered pair of classes of the two conjugate ideals is
\[
  (0,2),\quad(1,1),\quad\text{or}\quad(2,0).
\]
Theorem~\ref{thm:bad-class} says that class $2$ contributes no atom and
classes $0,1$ contribute.  Lemma~\ref{lem:one-atom} says that each
contributing ideal contributes exactly one atom.  Thus in the first
and third cases exactly one atom occurs, so the magnitude is $1$.  In
the middle case two atoms occur, and Lemma~\ref{lem:sign-coherence}
shows that they have the same sign, so the magnitude is $2$.  The
middle case is exactly $i=1$.  In all cases at least one atom occurs
and no cancellation to zero is possible.
\end{proof}

%----------------------------------------------------------------------
\section{Density from Hecke--Mitsui equidistribution}
\label{sec:density}
%----------------------------------------------------------------------

We now prove the analytic part of Theorem~\ref{thm:main}.  We keep
prime ideals and prime elements separate, since the Hecke--Mitsui
theorem naturally counts prime elements in an archimedean region with
a simultaneous finite congruence condition.

Recall that
\[
  \OK^\times=\{\pm\ph^n:n\in\ZZ\},
  \qquad \eps=\ph^2,
\]
that the norm-one units are $\{\pm\eps^m:m\in\ZZ\}$, and that the
totally positive units are $\{\eps^m:m\in\ZZ\}$.  For every nonzero
prime ideal $\mathfrak p$ not above $2$ or $5$, let
$\pi_{\mathfrak p}$ be the unique generator satisfying
\[
  \N(\pi_{\mathfrak p})=-\N(\mathfrak p),
  \qquad
  \sigma_1(\pi_{\mathfrak p})>0>
  \sigma_2(\pi_{\mathfrak p}),
  \qquad
  1\leq R(\pi_{\mathfrak p})<\rho.
\]
Class number one first gives a generator.  Multiplication by $\ph$ if
necessary makes its norm negative, and multiplication by $-1$ fixes
the sign pattern.  The remaining associates preserving both choices
are exactly $\eps^m\pi_{\mathfrak p}$.  Since
$R(\eps^m\alpha)=\rho^mR(\alpha)$, the half-open interval selects a
unique one.  If a prime element lies on a boundary ray, unit-normalizing
it as in the proof of Lemma~\ref{lem:no-boundary} gives
$\gamma=m\sqrt5$ with $m\in\ZZ$.  Since $\gamma$ is still a prime
element, primality forces $|m|=1$, so its principal ideal is the unique
prime above $5$.  Thus away from that prime we may use the open
inequalities $1<R<\rho$.

Put
\[
  \mathfrak m=(2\sqrt5)=2\sqrt5\,\OK.
\]
The ideals $(2)$ and $(\sqrt5)$ are coprime.  Since $2$ is inert and
$5$ is ramified, the Chinese remainder theorem gives
\[
  (\OK/\mathfrak m)^\times
  \simeq (\OK/2\OK)^\times\times
         (\OK/\sqrt5\OK)^\times
  \simeq \FF_4^\times\times\FF_5^\times,
  \qquad
  \#(\OK/\mathfrak m)^\times=12.
\]
This cardinality alone is not a ray-class argument.  Indeed, the
reduction of $\ph$ has order $3$ modulo $2$ and order $4$ modulo
$\sqrt5$, hence generates all twelve reduced classes modulo
$\mathfrak m$.  The equal main terms below come from the joint
archimedean-sector and finite-congruence theorem, where the unit group
acts diagonally; the canonical sector fixes the unit ambiguity.
For $\ell\in\ZZ/3\ZZ$ and $s\in\FF_5^\times$, let
$u_{\ell,s}\in(\OK/\mathfrak m)^\times$ be the unique class with
\[
  u_{\ell,s}\equiv\eps^\ell\pmod{2\OK},
  \qquad
  u_{\ell,s}\equiv s\pmod{\sqrt5\OK}.
\]
Let $\widetilde\Pi_{\ell,s}(X)$ count prime ideals
$\mathfrak p\nmid\mathfrak m$ of norm at most $X$ whose canonical
generator satisfies
$\pi_{\mathfrak p}\equiv u_{\ell,s}\pmod{\mathfrak m}$.

\begin{theorem}[Hecke--Mitsui input, specialized]
\label{thm:mitsui-special}
For every $\ell\in\ZZ/3\ZZ$ and every $s\in\FF_5^\times$,
\[
  \widetilde\Pi_{\ell,s}(X)
  =\frac1{12}\operatorname{Li}(X)
   +o\!\left(\frac{X}{\log X}\right).
\]
The same asymptotic holds after restricting to degree-one prime
ideals.
\end{theorem}

\begin{proof}
Mitsui's generalized prime number theorem gives simultaneous
archimedean and reduced-congruence equidistribution of prime elements;
see \cite{jMi56}, Corollary and Theorem on p.~35.  For the precise
modern form used here, see Kai's thin-cone theorem
\cite[Corollary~4.2.3]{oKai22}.  Apply it with the fixed modulus
$\mathfrak m$ and the open fundamental sector
\[
  \sigma_1(\pi)>0>\sigma_2(\pi),
  \qquad 1<R(\pi)<\rho.
\]
Choose once and for all a lower radial cutoff large enough to satisfy
the hypotheses of Kai's corollary.  The finitely many prime elements
below that cutoff contribute $O(1)$ and do not affect the asymptotic.
Partition its projectivized interval into finitely many sufficiently
small open thin-cone charts, choosing the internal cut rays to contain
no nonzero lattice point, and sum.  The two outer boundary rays contain
no prime element away from the prime above $5$, by the preceding
normalization.

For each reduced class $u\in(\OK/\mathfrak m)^\times$, the resulting
logarithmically weighted count has the form
\[
  \Theta_u(X)=C X+o(X),
\]
where $C$ is independent of $u$.  Kai's refinement includes a
possible exceptional-zero secondary term; because the modulus is
fixed, it is $O(X^\beta)$ for a fixed $\beta<1$ and hence is $o(X)$.
The canonical-generator normalization identifies the union over all
twelve reduced classes with the prime ideals away from
$\mathfrak m$.  Summing over those classes and applying the prime
ideal theorem gives $12C=1$.  Thus $C=1/12$, and partial summation
gives the displayed asymptotic.

Finally, in a quadratic field every prime ideal of residue degree two
and norm at most $X$ lies over a rational prime at most $\sqrt X$.
There are $O(\sqrt X/\log X)$ such ideals, so deleting them does not
change the main term.
\end{proof}

For a degree-one prime ideal $\mathfrak p$ of norm $p$, the canonical
generator has $\N(\pi_{\mathfrak p})=-p$.  Conjugation is trivial
modulo $\sqrt5$, so if
$\pi_{\mathfrak p}\equiv s\pmod{\sqrt5}$, then
\[
  -p\equiv\N(\pi_{\mathfrak p})
      \equiv\pi_{\mathfrak p}^{2}
      \equiv s^2\pmod5.
\]
Consequently the degree-one prime ideals above rational primes
$p\equiv1\pmod{10}$ are precisely the six selected cells
\[
  \ell\in\ZZ/3\ZZ,
  \qquad s\in\{2,3\}.
\]
The nonzero mod-$2$ component makes the norm odd, while $s=2,3$ is
equivalent to $s^2\equiv-1\pmod5$.  Conversely, both prime ideals
above every $p\equiv1\pmod{10}$ have canonical generators in these
six cells.

\begin{proposition}[Equal distribution of $\iota$ on prime ideals]
\label{prop:iota-density}
Let $I_i(X)$ count degree-one prime ideals of norm at most $X$ above
rational primes $p\equiv1\pmod{10}$ and satisfying
$\iota=i\in\ZZ/3\ZZ$.  Then
\[
  I_i(X)=\frac16\operatorname{Li}(X)
  +o\!\left(\frac{X}{\log X}\right)
  \qquad(i=0,1,2).
\]
In particular each value of $\iota$ has relative density $1/3$ among
these prime ideals.
\end{proposition}

\begin{proof}
For the canonical generator $\pi_{\mathfrak p}$, the sector index is
$0$, so
\[
  \iota(\mathfrak p)=-\lambda(\pi_{\mathfrak p})\pmod3.
\]
For each fixed $\ell$ there are exactly two selected cells, namely
$s=2$ and $s=3$.  Theorem~\ref{thm:mitsui-special} gives
$\frac1{12}\operatorname{Li}(X)$ for each cell, hence
$\frac16\operatorname{Li}(X)$ for each value of $\iota=-\ell$.
\end{proof}

\begin{theorem}[Rational-prime densities]\label{thm:density}
Let $P_2(X)$ and $P_1(X)$ count rational primes $p\leq X$ with
$p\equiv1\pmod{10}$ for which respectively
\[
  \left|B_{(p-1)/10}\right|=2
  \qquad\text{and}\qquad
  \left|B_{(p-1)/10}\right|=1.
\]
Then
\[
  P_2(X)=\frac1{12}\operatorname{Li}(X)
  +o\!\left(\frac{X}{\log X}\right),
\]
and
\[
  P_1(X)=\frac16\operatorname{Li}(X)
  +o\!\left(\frac{X}{\log X}\right).
\]
Consequently, relative to rational primes $p\equiv1\pmod{10}$,
\[
  \operatorname{dens}\left(
    \left|B_{(p-1)/10}\right|=2\right)=\frac13,
  \qquad
  \operatorname{dens}\left(
    \left|B_{(p-1)/10}\right|=1\right)=\frac23.
\]
\end{theorem}

\begin{proof}
Every such rational prime has exactly two conjugate prime ideals.
The reflection map on classes is
\[
  i\longmapsto-i-1,
\]
so it pairs $0$ with $2$ and fixes $1$.  By
Theorem~\ref{thm:prime-classification}, magnitude $2$ occurs exactly
for the conjugate class pair $(1,1)$, whereas magnitude $1$ occurs for
$(0,2)$ or $(2,0)$.  Consequently
\[
  I_1(X)=2P_2(X).
\]
Proposition~\ref{prop:iota-density} gives
\[
  P_2(X)=\frac1{12}\operatorname{Li}(X)
  +o\!\left(\frac{X}{\log X}\right).
\]
A magnitude-$1$ prime has one class-$0$ and one class-$2$ ideal, so
\[
  I_0(X)=I_2(X)=P_1(X),
\]
which gives the second asymptotic.  Equivalently,
\[
  \#\{p\leq X:p\equiv1\pmod{10}\}
  =\frac14\operatorname{Li}(X)
   +o\!\left(\frac{X}{\log X}\right).
\]
Dividing the two magnitude counts by this total gives the relative
densities $1/3$ and $2/3$.
\end{proof}

Theorems~\ref{thm:prime-classification} and \ref{thm:density} prove
Theorem~\ref{thm:main}.

%----------------------------------------------------------------------
\section{Computational checks and Lean scope}\label{sec:verification}
%----------------------------------------------------------------------

\subsection{Computational checks}

A direct termwise comparison of the two defining cones, signs, and
exponents gives the exact identification, in Hickerson--Mortenson
notation \cite{jHiMo14},
$B=-f_{1,3,4}(X,-X^3,X)$.  The coefficient implementation was also
checked through $N=200$.  The following statements are computational
checks, not additional inputs to the proofs above.  The aggregate cone and
prime statistics below---the $19{,}617$ prime sweep, the $7{,}165$ pair
sweep, the maximum coefficient, and the norm-supported zero count---are
reproduced by the repository script
\texttt{paper/numerical\_audit.py}.  The
three separate termwise, row-model, and ideal-level checks use their stated
finite ranges.

\begin{itemize}
\item The motivating row-model factorization has been verified
coefficient-by-coefficient through $N=2048$.  This computation is not
used as a source-faithful identification with Chan's missing kernel.
\item The prime classification has been checked for all $19{,}617$
primes $p\equiv1\pmod{10}$ with
$p\leq1{,}000{,}001$ (equivalently $(p-1)/10\leq100{,}000$).
No zero occurs in that exact range.
\item The explicit values at $11$, $31$, and $341$ are proved in
Theorem~\ref{thm:nonmult}; they are not being used merely as numerical
evidence.
\item The one-atom and sign-coherence lemmas were additionally
verified by direct ideal-level enumeration for every prime
$p\equiv1\pmod{10}$ with $p<20{,}000$ (and sampled up to $10^6$):
each contributing prime ideal received exactly one cone atom, and
paired contributions always carried equal signs.  In particular the
cases $p=11,31,41,61,71,101$ can be checked by hand.
\item Nonmultiplicativity is systematic: among the coprime pairs of
distinct primes $p<q$, $p,q\equiv1\pmod{10}$, with
$pq\leq10^6+1$ ($7165$ pairs),
none satisfies $a(pq)=a(p)a(q)$.
\item Through $N\leq10^5$ one has $\max_N|B_N|=7$, and there are
$4904$ levels in that range with at least one defining cone atom but
$B_N=0$.  Every corresponding value of $10N+1$ is composite; empirically
$|B_N|=O(d(10N+1))$.
\end{itemize}

\subsection{Lean scope}

The accompanying Lean~4 development has no admitted lemmas in the
algebraic, geometric, or finite prime-classification chain.  For every
rational prime $p\equiv1\pmod {10}$ it constructs an element of
$\ZZ[\phi]$ of norm $-p$, using quadratic reciprocity, a finite
pigeonhole argument, and a sign correction by the fundamental unit.
It classifies all norm-one units as $\pm\eps^m$, proves that the
norm-$-p$ elements lie in the two signed fundamental-unit orbits
represented by a generator and its conjugate, and constructs the unique
integral sector normalization.  It also proves that the resulting
generator-level $\iota$ value is unchanged when the norm-$-p$ generator is
replaced by another element in the same explicitly defined signed
$\eps$-power orbit.  The geometric
reflection is formalized with the exact transformed lift
$-\ell-1$.

The Lean normalization uses the integral trace/window formulation.  Its
equivalence with the real-embedding sector inequalities is the paper-level
content of Proposition~\ref{prop:cone-recognition}, not a separate Lean
endpoint.

The finite layer constructs a single finite type containing all cone
atoms at level $(p-1)/10$.  It proves both directions of the bridge
between the two signed fundamental-unit orbit classes and the contribution
predicate, proves that the two conjugate orbit classes are distinct (the
exceptional ramified prime $5$ being excluded by
$p\equiv1\pmod {10}$), and obtains the exact indicator formula for the atom
count.  The reflected classes contribute simultaneously exactly when the
geometric generator-level $\iota$ value is $1$.  A coordinate-level
conjugation theorem then proves equality of the full signed atom weights,
not merely of their signs.  Consequently Lean proves, from only primality
and $p\equiv1\pmod {10}$,
\[
 B_{(p-1)/10}\ne0,\qquad
 B_{(p-1)/10}\in\{-2,-1,1,2\}.
\]
For every norm-$-p$ generator $x$ and every geometric certificate $c$,
it also proves the exact magnitude-two equivalence expressed by
\texttt{prime\_\allowbreak BCoeff\_\allowbreak natAbs\_\allowbreak
eq\_\allowbreak two\_\allowbreak iff\_\allowbreak iota\_\allowbreak
one\_\allowbreak unconditional}.
The paper-level ideal interpretation of that generator invariant is the
well-definedness statement of Proposition~\ref{prop:iota-wd}.
The development also checks the explicit nonmultiplicativity witness
$B_1=1$, $B_3=-2$, $B_{34}=3$ and the theorem
\texttt{BCoeff\_\allowbreak not\_\allowbreak multiplicative\_\allowbreak
witness}.

The axiom audit for these endpoints reports only Lean's standard
\texttt{propext}, \texttt{Classical.\allowbreak choice}, and
\texttt{Quot.\allowbreak sound};
there are no project axioms or native-evaluation axioms.  The
specialized Hecke--Mitsui prime-element theorem used in
Section~\ref{sec:density} is not presently available in Mathlib and
remains an explicit analytic hypothesis.  A separate module in the broader
internal Lean development, not included in this public package, is a
schematic consumer: given $1/12$ asymptotics for the six selected cells
among the twelve reduced classes, together with the two exact equalities
identifying the relevant ideal counts with the magnitude-one and
magnitude-two rational-prime counts, it proves the resulting $1/6$,
$1/12$, $1/3$, and $2/3$ constant arithmetic.  It does not define the
concrete prime-ideal counting functions or prove those interface hypotheses.
A
source-faithful bridge from Chan's original missing kernel to the
independently defined series $B$ is not formalized and is not used in
the results stated here.

%----------------------------------------------------------------------
\section{Concluding remarks}\label{sec:conclusion}
%----------------------------------------------------------------------

The series $B$ shows that a norm-supported real-quadratic cone sum can
retain rigid prime behavior after the standard multiplicative model
has failed.  The obstruction to the usual descent is visible in the
unit action on the affine coset: $\eps L$ is disjoint from $L$, and
only $\eps^6$ stabilizes the full coset.  The explicit
$11\cdot31$ calculation proves nonmultiplicativity, while the unit
calculation explains why the classical one-coset construction is
unavailable.

At primes, the replacement for the unavailable one-coset Hecke
character model is the
sector--residue invariant $\iota$.  Its reflection law gives the three
possible conjugate class pairs.  The one-atom lemma turns each
good ideal into exactly one lattice point, and sign coherence prevents
the pair in class $(1,1)$ from cancelling.  This yields nonvanishing
and the two possible magnitudes.  The specialized joint
prime-element theorem gives the exact density of each of the twelve
sector--congruence cells.  The six cells over primes
$p\equiv1\pmod{10}$ are therefore equidistributed, and explicit
two-ideal bookkeeping transfers the resulting counts to rational
primes.

\bibliographystyle{plain}
\bibliography{master}

\end{document}
